\vfill\pagebreak
\section{Integer types}
\label{sec:types:integers}

Ymir has ten integer types, and their names say what they are. The letter is the
sign: \token{i8}, \token{i16}, \token{i32}, \token{i64} and \token{isize} are
signed and hold negative values as well as positive ones; \token{u8},
\token{u16}, \token{u32}, \token{u64} and \token{usize} are unsigned and hold
only non-negative ones. The number is the width in bits, and the width fixes how
many distinct values the type can tell apart. \token{isize} and \token{usize}
are the exception: their width is whatever the machine works in naturally, 64
bits on ARM64, 32 on x86.

\subsection{Literal representation}
\label{sec:types:integer_literals}

An integer literal is a run of digits written directly in the source code, in
decimal and without spaces -- \token{12} and \token{36} in the listings so far.
Long ones are hard to read, so underscores may be scattered through them as
separators: \token{1\_000\_000\_000\_000} is one trillion, and the underscores
are ignored.

A literal with nothing else attached is an \token{i32}. A suffix names another
type: \token{12u8} is an 8-bit unsigned value, \token{12i64} a 64-bit signed
one.

\notebox{
  More advanced integer literals can be written in octal, hexadecimal,
  or binary form. These representations are covered in detail in the advanced
  chapter on integers (see Section~\ref{sec:scalars:integers}).
}

\subsection{Memory representation}
\label{sec:types:integer_memory}

A computer stores numbers in base~$2$. Each digit is a bit, worth $0$ or $1$, and
each position is worth twice the one to its right: $2^0$ at the far right, then
$2^1$, $2^2$, and so on. An 8-bit integer is thus an eight-digit number in
base~$2$, and occupies one byte.

\begin{table}[h!]
  \centering
  \label{tab:types:integer_ranges}
  %% Columns sized to their contents rather than to fixed fractions of the text
  %% width. As p{} columns, the Type cells were too narrow for their two chips
  %% and the Signed Range cells too narrow for their formula, so both wrapped --
  %% which put each type one line below the row it labels, and broke $2^{N-1}-1$
  %% across two lines after the minus sign.
  \begin{tabular}{l|l|c|c}
    \toprule[0.6pt]
    \textbf{Type} & \textbf{Bits} & \textbf{Signed Range} & \textbf{Unsigned Range} \\
    \toprule[0.6pt]
    \token{i8} / \token{u8}      & 8    & $-2^7 \ldots 2^7-1$        & $0 \ldots 2^8-1$ \\
    \token{i16} / \token{u16}    & 16   & $-2^{15} \ldots 2^{15}-1$  & $0 \ldots 2^{16}-1$ \\
    \token{i32} / \token{u32}    & 32   & $-2^{31} \ldots 2^{31}-1$  & $0 \ldots 2^{32}-1$ \\
    \token{i64} / \token{u64}    & 64   & $-2^{63} \ldots 2^{63}-1$  & $0 \ldots 2^{64}-1$ \\
    \token{isize} / \token{usize} & system-dependent (N) & $-2^{N-1} \ldots 2^{N-1}-1$ & $0 \ldots 2^N-1$ \\
    \midrule[0.2pt]
  \end{tabular}
  \caption{Integer Types and Ranges in Ymir (expressed as powers of 2)}
\end{table}



Unsigned 8-bit integers \token{u8} use all 8 bits to represent non-negative
numbers, ranging from $0$ to $2^8-1 = 255$. Each bit contributes to the total
value depending on its position: the rightmost bit represents 1, the next bit 2,
then 4, 8, 16, 32, 64, and finally 128 for the leftmost bit. For instance, the
number 13 in binary is $00001101$, which corresponds to:
\begin{equation}
  (0\times{}128) +
  (0\times{}64) + (0\times{}32) + (0\times{}16) + (1\times{}8) + (1\times{}4) + (0\times{}2) +
  (1\times{}1) = 13
\end{equation}

A signed \token{i8} has to fit negative numbers into the same eight bits, and
does it by two's complement. The leftmost bit -- the most significant, or MSB --
carries the sign: $0$ for positive, $1$ for negative. Positive values are
unchanged, so $13$ is still $00001101$. A negative value is built in three steps:
write the magnitude in binary ($00001101$), invert every bit ($11110010$), add
one ($11110011$). That last pattern is $-13$.

Spending a bit on the sign costs range but not capacity. An \token{i8} still
distinguishes 256 values, exactly as a \token{u8} does; they are simply a
different 256, from $-2^7$ to $2^7-1$, which is $-128$ to $127$.

\notebox{
  The unary operator \token{!}, when applied to an integer, returns its one’s
  complement—that is, it inverts all the bits of the number. For example, the
  operation \token{(!13) + 1} produces the two’s complement of $13$, which is equal
  to \token{-13}.
}

Nothing above depends on the width being eight. For $N$ bits, an unsigned type
covers $0$ to $2^N-1$ and a signed one $-2^{N-1}$ to $2^{N-1}-1$. Base~$2$ can
write any whole number given enough digits; a machine has a fixed number of them,
which is why every integer type has ends.


\subsection{Integer overflow}
\label{sec:types:integer_overflow}

Those ends are hard. When a calculation lands outside them, the result is not the
number that was wanted -- it is \textit{integer overflow}, and the value wraps
around to the other end of the range.

A \token{u8} runs from $0$ to $255$, so $255 + 1$ is $0$ and $0 - 1$ is $255$. An
\token{i8} runs from $-128$ to $127$, so $127 + 1$ is $-128$. Ymir does not
detect this and does not prevent it. It cannot be caught at compile time, since
the operands are generally not known until the program runs, and checking every
arithmetic operation at run time would cost more than it is worth.

The consequence is worth stating plainly: an overflowed calculation produces a
wrong answer silently, and the program carries on with it. Knowing the range of
the type in hand is part of choosing it.
